\section{实验设置与结果分析}
\label{rq3:sec:experiments}

\subsection{任务数据与实验流程}
\label{rq3:subsec:experiment-data}

实验采用基于 VQAv2/COCO 图像构造的冻结视觉语义轨迹\citep{coco,vqa-v2}。每个任务由一幅图像及其查询构成，通过局部裁剪为九个感知节点提供不同观测，接收中心不直接访问原图。离线视觉语言模型生成局部证据和全图参考标签，标签空间为清晰、基本清晰、不确定、基本遮挡和遮挡五个等级。因此，本章的任务正确性表示证据可见性分类与全图模型参考标签一致，评价对象是给定冻结语义内容后的通信效果。

候选证据由前期选择性语义通信模块 EviComRL 生成。该模块采用包含结构化候选特征与历史上下文的 Actor-Critic 策略，并以 PPO 进行训练\citep{ppo}。本章使用其在本实验环境训练后固定的检查点，在预设 token 预算下生成候选序列，经固定顺序合并、去重后保留至多四条发往中心的证据。所有比较方法使用相同候选内容与次序。整个反馈控制实验不重新调用视觉语言模型，也不联合微调候选选择模块。

语义效用学习使用训练分区中的 600 个任务，验证分区包含 595 个任务，测试分区包含 611 个任务。图像在不同分区之间不重叠。每个训练任务枚举所有非空接收集合，形成式~\eqref{rq3:eq:utility-loss} 的监督样本；同一图像产生的不同集合仍属于同一个原始任务。效用网络采用 AdamW 优化，学习率为 $8\times10^{-4}$，权重衰减为 0.02，批量大小为 256，最多训练 60 轮，并依据验证交叉熵选择参数。

控制参数的选择在四种验证条件下进行，分别为参考条件、12 ms 反馈延迟、60\% 反馈丢失率和 12 ms 任务时限。对于每种方法及训练种子，首先比较四种条件的平均正确完成率，再以更低通信能耗处理并列情况，最终保留一组参数用于全部测试条件。C-BFC 的三个训练种子均选择 $\eta_E=0.24$、$\eta_T=0$。时间约束仍通过发送可行性、剩余规划时域和等待机会成本影响动作选择。

包含学习模型的方法使用三个训练种子与三个网络种子的九种组合，固定协议使用三个网络种子。每次测试按相同任务次序处理 611 个任务，节点能量在该次任务流内连续扣减。不同方法对同一任务、同一候选和同一次尝试使用相同随机条件，以减少抽样差异对方法比较的影响。图中误差线表示重复实验的标准差。重复实验改变模型初始化和网络随机条件，视觉测试样本保持固定。

\subsection{通信参数与评价指标}
\label{rq3:subsec:experiment-settings}

表~\ref{rq3:tab:parameters} 给出了主要参数。节点到中心的距离在 100--260 m 范围内平滑变化，可用带宽系数在标称带宽的 0.25--1.0 之间变化，分别模拟任务流中传播条件与共享无线资源的波动。速率预测按式~\eqref{rq3:eq:rate} 计算，单次分组成功按式~\eqref{rq3:eq:success-probability} 抽样。每个任务最多执行九次数据发送，不同候选可能来自同一发送节点，因此能量约束始终按节点汇总。

\begin{table}[t]
  \centering\small
  \bicaption{语义交付实验的主要通信与控制参数}{Main communication and control parameters for semantic delivery}
  \label{rq3:tab:parameters}
  \begin{tabularx}{\linewidth}{@{}Yl@{}}
    \toprule
    参数 & 参考值或设置\\
    \midrule
    感知节点数／接收中心数 & 9／1\\
    每任务候选上限 $K_{\max}$ & 4\\
    标称带宽 & 20 kHz\\
    数据／反馈发射功率 & 0.2 W／0.1 W\\
    参考信道增益 $\rho_0$ & $10^{-4}$\\
    噪声功率谱密度 $N_0$ & $4\times10^{-20}$ W/Hz\\
    每 token 字节数 $b$ & 4\\
    数据头部／有效载荷上限 & 64 B／1200 B\\
    反馈帧长度 & 24 B\\
    发送时长裕量 $\kappa$ & 1.6\\
    单次发送时长上限 $\tau_{\max}$ & 8 ms\\
    任务时限 $D_k$ & 25 ms\\
    反馈平均额外时延 $\bar d_{\mathrm f}$ & 6 ms\\
    反馈丢失率 $\epsilon_{\mathrm f}$ & 25\%\\
    速率波动强度 $\sigma_r$ & 0.55\\
    每节点初始发送能量 $E_i^{(0)}$ & 0.15 J\\
    数据尝试上限／基础等待间隔 & 9 次／2 ms\\
    C-BFC 能量／时间权重 & 0.24／0\\
    \bottomrule
  \end{tabularx}
\end{table}

实验使用四项指标共同评价通信效果。

\paragraph{按时正确完成率}
任务必须在截止时间内至少接收一条证据，且最终类别与参考标签一致，才计为正确完成。由于环境仅执行可在截止期内完整结束的数据动作，该指标为
\begin{equation}
  R_{\mathrm c}=\frac{1}{N_{\mathrm{task}}}\sum_{k=1}^{N_{\mathrm{task}}} u_k(S_k^{\mathrm{end}}).
  \label{rq3:eq:metric-correct}
\end{equation}
分母包含因能量不足或策略停止而未获得服务的任务。

\paragraph{单位正确任务通信能耗}
将整个任务流的数据与反馈发送能量除以正确完成的任务数：
\begin{equation}
  E_{\mathrm c}=\frac{\sum_{k=1}^{N_{\mathrm{task}}} E_k^{\mathrm{comm}}}
                         {\sum_{k=1}^{N_{\mathrm{task}}} u_k(S_k^{\mathrm{end}})}.
  \label{rq3:eq:metric-energy}
\end{equation}
该指标以 mJ 表示，可同时反映通信支出与正确交付效果。先对每次回放计算式~\eqref{rq3:eq:metric-energy}，再跨重复实验与条件求平均。

\paragraph{重复发送比例}
若某次发送开始前，同一证据已成功到达接收端，则该次尝试属于重复发送。令 $d_{k,a}$ 为相应指示变量，重复比例为
\begin{equation}
  R_{\mathrm{dup}}=\frac{\sum_k\sum_{a\in\mathcal T_k}d_{k,a}}
                         {\sum_k|\mathcal T_k|}.
  \label{rq3:eq:metric-duplicate}
\end{equation}
这一指标依据实际接收日志计算，仅用于评价控制器未能避免的冗余。已尝试但此前从未成功的补发不计为重复发送。

\paragraph{平均通信完成时延}
记 $\mathcal K_+=\{k:|\mathcal T_k|>0\}$ 为实际发起数据发送的任务集合，则
\begin{equation}
  T_{\mathrm{comm}}=\frac{1}{|\mathcal K_+|}\sum_{k\in\mathcal K_+}T_k.
  \label{rq3:eq:metric-latency}
\end{equation}
该指标包含任务中的发送与等待时间，但不把直接跳过的任务以零时延计入平均。较低时延可能来自及时停止，也可能来自发送较少，因此需要结合正确完成率解释。

\subsection{对比方法与实验分组}
\label{rq3:subsec:baselines}

对比方法包括固定冗余、停等 ARQ、选择重传、语义阈值 ARQ、队列贪心，以及两个学习型控制器。固定冗余按照验证选择的候选数量和副本数发送，不依赖反馈触发后续动作。停等 ARQ 等待累计确认或超时后继续发送；选择重传允许有限窗口内的未确认候选独立补发。语义阈值 ARQ 在累计确认对应的语义效用达到阈值时结束任务。队列贪心使用接收信念与能量价格，逐步选择即时净收益最高的发送动作。

SD-SAC 对照保留文献中的双 Q 平均、Q 值裁剪与熵变化惩罚\citep{zhou2024sdsac}。TP-DQN 参考两阶段码率控制的思路，先使用语义性能预测模型给出任务效用，再训练离散动作价值网络\citep{wu2026twophase}，其价值更新采用 Double DQN 形式。两者均适配为本章的反馈动作空间：四个候选发送动作、三种等待动作和停止动作。策略输入包含候选属性、接收状态概率、语义效用与资源条件，共 89 维，使用 $89\rightarrow128\rightarrow64\rightarrow8$ 的网络结构。

两个学习型控制器都在当前环境重新训练，每个训练种子执行 80 个由 600 个任务构成的任务流，并在十二种物理与反馈设置之间变化。训练每四次交互执行一次更新，批量大小为 128，学习率为 $3\times10^{-4}$。SD-SAC 使用训练任务的正确完成奖励，TP-DQN 使用学习得到的语义性能评价；二者均加入效用势函数差作为训练信号。最终测试使用验证选中的在线训练参数。这里比较的是相同语义交付任务中的控制机制，原研究中的文本编解码网络不属于本实验环境。

实验按五组条件展开：速率波动强度取 0.20、0.35、0.55、0.80、1.05；反馈平均额外时延取 0、3、6、9、12 ms；任务时限取 10、15、25、35、45 ms；每节点发送能量取 0.06、0.10、0.15、0.24、0.36 J；组件实验在 0\%、15\%、25\%、45\%、60\% 五档反馈丢失率下进行。每组只改变对应参数，其他参数保持参考值。前四组去除重复参考点后共有十七个主比较条件，连同组件实验共二十一个不同条件，总计完成 1,089 次完整任务流回放。

主图展示固定冗余、语义阈值 ARQ、队列贪心、SD-SAC、TP-DQN 和 C-BFC，性能表同时列出停等 ARQ 与选择重传结果。含语义估计的对照使用相同训练种子的原始效用模型，C-BFC 在此基础上执行一致性投影。各方法的阈值、重传规则或资源权重均在验证分区选择后固定。

\subsection{总体性能比较}
\label{rq3:subsec:overall-results}

表~\ref{rq3:tab:main-results} 给出了十七个主条件的等权平均结果。C-BFC 的正确完成率为 49.84\%，单位正确任务通信能耗为 3.615 mJ。与固定冗余和语义阈值 ARQ 相比，C-BFC 分别提高正确完成率 9.35 和 6.45 个百分点。固定发送副本能够提高部分证据的交付机会，但无法根据已发生的传输结果及时调整；语义阈值 ARQ 能够利用已确认集合的效用停止发送，但在反馈未返回时仍主要受超时规则支配。

\begin{table}[t]
  \centering\small
  \setlength{\tabcolsep}{3.5pt}
  \bicaption{十七个主比较条件下的平均性能}{Average performance across the seventeen main evaluation conditions}
  \label{rq3:tab:main-results}
  \begin{tabular*}{\linewidth}{@{\extracolsep{\fill}}lrrrr@{}}
    \toprule
    方法 & \shortstack{正确完成率\\(\%)} & \shortstack{单位正确任务\\能耗 (mJ)} & \shortstack{重复发送\\(\%)} & \shortstack{通信时延\\(ms)}\\
    \midrule
    固定冗余 & 40.49 & 4.941 & 29.35 & 11.32\\
    停等 ARQ & 43.65 & 3.678 & 7.49 & 16.69\\
    选择重传 & 43.57 & 3.665 & 7.16 & 16.77\\
    语义阈值 ARQ & 43.39 & 4.090 & 6.71 & 17.63\\
    队列贪心 & 47.26 & 4.564 & 25.76 & 10.96\\
    SD-SAC & 48.97 & 4.334 & 15.67 & 17.76\\
    TP-DQN & 49.18 & 4.219 & 19.11 & 17.13\\
    \midrule
    \textbf{C-BFC} & \textbf{49.84} & \textbf{3.615} & \textbf{17.14} & \textbf{16.47}\\
    \bottomrule
  \end{tabular*}
  \par\vspace{4pt}{\footnotesize 各条件等权汇总，共享的参考条件只计入一次。加粗行表示本章方法。}
\end{table}


队列贪心的正确完成率为 47.26\%，说明显式接收信念和语义收益能够改善简单重传规则的任务效果。但是，其单位正确任务能耗为 4.564 mJ，重复发送比例为 25.76\%。这一结果与其仅评价当前发送收益、缺少主动等待相符：当接收状态尚不确定时，连续执行即时收益为正的动作，可能消耗本可通过等待确认而节省的能量。

相较于 SD-SAC 和 TP-DQN，C-BFC 的正确完成率分别提高 0.87 和 0.66 个百分点，单位正确任务能耗分别降低 16.59\% 和 14.32\%，通信完成时延分别降低 7.29\% 和 3.87\%。在十七个主条件中，C-BFC 相对于两者的正确完成率分别在十二个和十四个条件下更高，单位正确任务能耗在全部十七个条件下更低。其平均正确率优势较小，主要收益表现为在相近或更高正确交付水平下减少通信支出。

\begin{table}[t]
  \centering\small
  \setlength{\tabcolsep}{3.5pt}
  \bicaption{参考条件下的性能比较}{Performance comparison under the reference condition}
  \label{rq3:tab:reference-results}
  \begin{tabular*}{\linewidth}{@{\extracolsep{\fill}}lrrrr@{}}
    \toprule
    方法 & \shortstack{正确完成率\\(\%)} & \shortstack{单位正确任务\\能耗 (mJ)} & \shortstack{重复发送\\(\%)} & \shortstack{通信时延\\(ms)}\\
    \midrule
    固定冗余 & 41.68 & 5.023 & 29.86 & 11.79\\
    停等 ARQ & 45.66 & 3.627 & 7.72 & 17.46\\
    选择重传 & 45.50 & 3.615 & 7.40 & 17.46\\
    语义阈值 ARQ & 46.32 & 3.948 & 7.07 & 18.16\\
    队列贪心 & 47.66 & 4.706 & 27.20 & 11.36\\
    SD-SAC & 49.95 & 4.356 & 15.66 & 18.24\\
    TP-DQN & 49.86 & 4.329 & 21.49 & 17.74\\
    \midrule
    \textbf{C-BFC} & \textbf{51.08} & \textbf{3.639} & \textbf{17.40} & \textbf{17.56}\\
    \bottomrule
  \end{tabular*}
  \par\vspace{4pt}{\footnotesize 参考条件为 25 ms 任务时限、6 ms 反馈延迟、25\% 反馈丢失率及每节点 0.15 J 发送能量。}
\end{table}


参考条件下的结果见表~\ref{rq3:tab:reference-results}。C-BFC 的正确完成率为 51.08\%，高于两个学习型对照约 1.1 个百分点，单位正确任务能耗由约 4.3 mJ 降至 3.639 mJ。停等与选择重传具有较低重复比例，说明保守确认规则有助于避免对已接收证据再次发送；与此同时，其正确完成率约为 45.5\%，表明较少重复并不必然对应更好的任务完成效果。因而，本章将正确完成率与单位正确任务能耗作为主要结果，并通过重复比例和时延解释其形成过程。

\FloatBarrier
\subsection{信道波动对交付效果的影响}
\label{rq3:subsec:fading-results}

图~\ref{rq3:fig:fading} 比较不同速率波动强度下的性能。在预测平均速率不变时，较大的 $\sigma_r$ 使单次发送更容易无法在预留时长内完成，因而多数方法的正确完成率下降，单位正确任务能耗上升。C-BFC 的正确完成率由 $\sigma_r=0.20$ 时的 57.26\% 下降至 $\sigma_r=1.05$ 时的 41.32\%，说明有限重传与反馈控制仍然受数据链路可靠性的约束。

\begin{figure}[!htbp]
  \centering
  \resultfigure{01_fading.pdf}{不同速率波动强度下的性能比较}
  \bicaption{速率波动强度变化下的性能比较}{Performance under different rate fluctuation levels}
  \par\vspace{2pt}{\footnotesize （a）按时正确完成率；（b）单位正确任务通信能耗；（c）重复发送比例；（d）平均通信完成时延。其余参数取表~\ref{rq3:tab:parameters} 的参考值。}
  \label{rq3:fig:fading}
\end{figure}

信道较稳定时，C-BFC 的节能优势更为明显。在 $\sigma_r=0.20$ 时，其单位正确任务能耗为 2.346 mJ，SD-SAC 和 TP-DQN 分别为 3.464 和 2.831 mJ，而三个方法的正确完成率均在 57\% 左右。此时，多次重传的实际需求较低，对接收状态与分类结果一致性进行判断能够较早识别额外发送的有限收益。

当 $\sigma_r=1.05$ 时，C-BFC 的单位正确任务能耗增至 5.036 mJ，仍低于 SD-SAC 的 5.621 mJ 和 TP-DQN 的 5.610 mJ。与此同时，各方法之间的正确完成率差距缩小。该结果表明，控制策略可以改善有限资源的使用方式，但无法完全抵消低成功概率对语义交付造成的损失。

\FloatBarrier
\subsection{确认反馈延迟的影响}
\label{rq3:subsec:feedback-results}

图~\ref{rq3:fig:feedback} 给出了反馈平均额外时延变化下的结果。额外时延为零时，反馈仍需完成串行发送，因而并非瞬时获得确认。随着额外时延增加，发送端维持接收不确定性的时间变长，等待确认与继续发送之间的矛盾更加明显。

\begin{figure}[!htbp]
  \centering
  \resultfigure{02_feedback.pdf}{不同反馈延迟下的性能比较}
  \bicaption{反馈平均额外时延变化下的性能比较}{Performance under different additional feedback delays}
  \par\vspace{2pt}{\footnotesize 四个子图依次对应正确完成率、单位正确任务能耗、重复发送比例和通信完成时延。}
  \label{rq3:fig:feedback}
\end{figure}
\FloatBarrier

在额外时延为 0 ms 时，C-BFC 的正确完成率为 53.52\%，通信完成时延为 11.36 ms；SD-SAC 与 TP-DQN 的正确完成率分别为 50.95\% 和 52.01\%，通信完成时延分别为 15.60 和 14.42 ms。较及时的确认能够使控制器尽早确定已经到达的证据，并据此停止冗余发送或选择互补内容。

当反馈平均额外时延增至 12 ms，C-BFC 的正确完成率降至 48.92\%，单位正确任务能耗由 3.447 mJ 增至 3.973 mJ，重复发送比例也随之增加。此时部分确认无法在有用的决策窗口内返回，主动等待的作用受到限制。两个学习型对照在同一条件下的正确完成率分别为 48.10\% 和 48.52\%，单位正确任务能耗分别为 4.426 和 4.526 mJ。C-BFC 仍保持较低的单位正确任务能耗，但其相对收益没有随着反馈延迟无限扩大。

\FloatBarrier
\subsection{任务截止时间的影响}
\label{rq3:subsec:deadline-results}

图~\ref{rq3:fig:deadline} 比较不同任务时限下的性能。时限同时决定能够执行的数据尝试数量和可用于等待确认的时间。10 ms 条件下，许多任务只能够容纳一次主要发送，控制器很难通过“发送后等待、再依据反馈补充证据”的过程获得额外收益。C-BFC 与 SD-SAC 的正确完成率分别为 44.48\% 和 44.64\%，差异较小。

\begin{figure}[!htbp]
  \centering
  \resultfigure{03_deadline.pdf}{不同任务截止时间下的性能比较}
  \bicaption{任务截止时间变化下的性能比较}{Performance under different task deadlines}
  \par\vspace{2pt}{\footnotesize 较长任务窗口允许反馈到达后继续调整发送，各方法均受到跨任务节点电量限制。}
  \label{rq3:fig:deadline}
\end{figure}

随着时限增加，C-BFC 的正确完成率提升，同时单位正确任务能耗保持在较低水平。在 45 ms 条件下，其正确完成率达到 54.74\%，较 SD-SAC 的 51.34\% 提高 3.40 个百分点；单位正确任务能耗为 3.539 mJ，较 SD-SAC 的 4.466 mJ 降低 20.76\%。与 TP-DQN 相比，C-BFC 也在正确完成率和单位正确任务能耗上取得改善。

C-BFC 在 45 ms 任务窗口下的平均通信完成时延为 23.39 ms，低于 SD-SAC 的 28.20 ms 和 TP-DQN 的 26.03 ms。算法能够利用返回的确认调整证据发送，并在继续通信的收益不足时提前结束。结合紧时限结果可见，反馈控制的作用取决于任务窗口能否容纳信息返回后的后续动作。

\FloatBarrier
\subsection{节点发送能量的影响}
\label{rq3:subsec:energy-results}

图~\ref{rq3:fig:energy} 将每节点初始发送能量从 0.06 J 增加到 0.36 J。能量额度较低时，当前任务中的支出会较快影响后续任务的可执行范围，正确完成率受到整个任务流资源总量的限制。在每节点 0.06 J 条件下，C-BFC 的正确完成率为 36.32\%，单位正确任务能耗为 2.427 mJ，说明其通过限制低收益发送保持了较高的资源使用效率。

\begin{figure}[!htbp]
  \centering
  \resultfigure{04_energy.pdf}{不同节点发送能量下的性能比较}
  \bicaption{每节点初始发送能量变化下的性能比较}{Performance under different initial node transmission energies}
  \par\vspace{2pt}{\footnotesize 电量在同一次 611 个任务的回放中连续扣减，其他通信条件固定。}
  \label{rq3:fig:energy}
\end{figure}

随着能量增加，各方法能够为更多任务提供通信机会，也能够执行更多补发。C-BFC 在每节点 0.10 J 时的正确完成率为 46.03\%，高于 SD-SAC 的 44.66\%；在参考能量 0.15 J 下进一步达到 51.08\%。这一变化说明，虚拟能量压力与剩余额度能够使同一组控制权重随资源条件改变实际发送行为。

在较高能量额度下，需要同时观察质量与能耗。每节点 0.36 J 时，C-BFC 的正确完成率为 55.74\%，低于 SD-SAC 的 57.57\% 和 TP-DQN 的 56.76\%；对应单位正确任务能耗为 4.608 mJ，低于两者的 6.164 和 6.033 mJ。此时，C-BFC 保留的能量偏好使其放弃了一部分代价较高的交付机会。因而，高能量条件下的结果体现了正确率与通信消耗之间的取舍，不能仅根据节能幅度判断所有性能指标均有改善。

\FloatBarrier
\subsection{关键组件作用分析}
\label{rq3:subsec:ablation}

为分析主要组件的作用，构造三个对照版本：无接收信念版本只使用最新确认位图，不传播未确认尝试的成功概率，也不计算依赖该分布的等待信息价值；无主动等待版本保留发送前瞻，但不将等待作为可主动选择的动作；无一致性约束版本直接使用原始效用预测。各版本在验证分区分别选择资源权重，再在五档反馈丢失率下进行测试。图~\ref{rq3:fig:ablation} 展示逐条件变化，表~\ref{rq3:tab:ablation-results} 给出等权平均结果。

\begin{figure}[!htbp]
  \centering
  \resultfigure{05_ablation.pdf}{不同反馈丢失率下的组件比较}
  \bicaption{不同反馈丢失率下的组件比较}{Component comparison under different feedback loss rates}
  \par\vspace{2pt}{\footnotesize 完整方法与无接收信念、无主动等待、无一致性约束三个版本分别比较，子图指标与前四组一致。}
  \label{rq3:fig:ablation}
\end{figure}

\begin{table}[t]
  \centering\small
  \setlength{\tabcolsep}{3.5pt}
  \bicaption{五档反馈丢失率下的组件比较}{Component comparison at five feedback loss rates}
  \label{rq3:tab:ablation-results}
  \begin{tabular*}{\linewidth}{@{\extracolsep{\fill}}lrrrr@{}}
    \toprule
    方法 & \shortstack{正确完成率\\(\%)} & \shortstack{单位正确任务\\能耗 (mJ)} & \shortstack{重复发送\\(\%)} & \shortstack{通信时延\\(ms)}\\
    \midrule
    无接收信念 & 39.69 & 5.814 & 60.46 & 16.91\\
    无主动等待 & 48.89 & 4.516 & 32.59 & 10.96\\
    无一致性约束 & 51.28 & 4.064 & 16.96 & 18.63\\
    \midrule
    \textbf{C-BFC} & \textbf{50.97} & \textbf{3.674} & \textbf{18.29} & \textbf{17.62}\\
    \bottomrule
  \end{tabular*}
  \par\vspace{4pt}{\footnotesize 各版本分别完成验证参数选择；表中数值为五个反馈丢失条件的等权平均。}
\end{table}


去除接收信念后，平均正确完成率由完整方法的 50.97\% 降至 39.69\%，单位正确任务能耗由 3.674 mJ 增至 5.814 mJ，重复发送比例达到 60.46\%。只依赖已确认位图会使控制器在反馈缺失时持续低估接收端掌握的信息，进而反复向已接收证据分配资源。这一结果支持显式区分数据丢失与反馈缺失的必要性。

去除主动等待后，平均正确完成率为 48.89\%，单位正确任务能耗为 4.516 mJ，重复发送比例为 32.59\%。完整方法通过等待可能到达的确认，减少了状态不确定时的连续发送，但平均通信完成时延也由 10.96 ms 增至 17.62 ms。可见，等待的价值表现为用一部分可用时间换取更合适的后续发送选择，其作用不能仅通过时延这一指标衡量。

去除一致性约束后，平均正确完成率为 51.28\%，略高于完整方法的 50.97\%；单位正确任务能耗为 4.064 mJ，高于完整方法的 3.674 mJ。加入一致性约束使单位正确任务能耗降低约 9.6\%，并伴随 0.31 个百分点的正确率变化。该结果与一致性投影的设计目的相符：抑制不改变分类结果的估计效用增量，主要改善通信支出，而不要求每个条件的正确率都提高。

\FloatBarrier
\subsection{结果讨论}
\label{rq3:subsec:discussion}

综合五组实验，C-BFC 的作用主要体现在接收不确定性、反馈等待与语义冗余共同存在的条件下。接收信念减少对确认缺失的误判，等待价值决定何时值得为信息返回保留时间，一致性效用使资源分配与接收端分类结果保持联系。三个部分共同改善了单位正确任务的通信消耗，但其收益受到链路成功概率、任务窗口和剩余电量的约束。

本章实验也给出了该方法的适用条件。首先，任务语义来自冻结轨迹，能够保证不同控制器面对相同候选与融合规则，但没有覆盖视觉语言模型在线生成内容发生变化的情况。其次，接收集合依赖至多四条候选，状态枚举的开销受到这一规模限制。再次，数据与确认采用独立抽样的发送模型，尚未考虑强相关衰落、反馈冲突或软信息合并。因此，本章结论对应调度层的语义可靠交付，不直接等同于完整无线原型系统的性能。

从资源配置角度看，较短时限首先限制反馈的可用程度，较低能量则限制整个任务流中的服务机会；资源较充足时，控制器仍需在少量额外正确完成与较高通信支出之间选择。本章通过固定验证参数考察这种条件变化，后续可进一步研究面向不同服务质量要求的工作点选择，以适应可靠性与能耗偏好不同的应用。
